% Original: PDF strana 22, gornji slajd.
\slajdid{02-s043}
\begin{frame}[label=02-s043]{Primer: minimalni zbir proizvoda}
\setlength{\parskip}{0pt}
Na osnovu zahteva za kombinacionom mrežom formirana je polazna Karnoova tabela.
\par\vspace{1.5mm}
\begin{columns}[T,onlytextwidth]
\begin{column}{.29\textwidth}\centering
% element: 02-s043:e01
\Oznaka{Polazna karta}\par\vspace{2mm}
% Bitovi indeksa su DCBA; kolone BA=00,01,11,10, redovi DC=00,01,11,10.
\begingroup
\fontsize{16.67}{18.5}\selectfont
\renewcommand{\small}{\fontsize{16.67}{18.5}\selectfont}
\scalebox{.54}{%
\begin{karnaugh-map}(label=middle)[4][4][1][$A$][$B$][$C$][$D$]
\minterms{0,4,5,8,10}
\terms{1}{b}
\maxterms{2,3,6,7,9,11,12,13,14,15}
\end{karnaugh-map}
}
\endgroup

\end{column}
\begin{column}{.29\textwidth}\centering
% element: 02-s043:e03
\Oznaka{Pokrivanje jedinica}\par\vspace{2mm}
\begingroup
\fontsize{16.67}{18.5}\selectfont
\renewcommand{\small}{\fontsize{16.67}{18.5}\selectfont}
\scalebox{.54}{%
\begin{karnaugh-map}(label=middle,implicantcolors={DEcrvena,DEcrvena,DEcrvena})[4][4][1][$A$][$B$][$C$][$D$]
\minterms{0,4,5,8,10}
\terms{1}{b}
\maxterms{2,3,6,7,9,11,12,13,14,15}
\implicant{0}{5}
{\tikzset{every path/.append style={dashed}}\implicantedge{8}{8}{10}{10}}
\end{karnaugh-map}
}
\endgroup

\end{column}
\begin{column}{.38\textwidth}
% element: 02-s043:e02
Za minimalnu funkciju u obliku zbira proizvoda prekrivamo sve logičke jedinice površinama najvišeg reda, koristeći po potrebi i neodređeno stanje \(b\).
\par\vspace{1.5mm}
Površine objedinjuju proizvode koji se mogu sažeti. U proizvodima ostaju samo literali koji se \alert{\textbf{ne menjaju}} na površini.
\end{column}
\end{columns}
\par\vspace{2mm}
% element: 02-s043:e04
\par\Rezultat{\(F=\bar D\bar B+D\bar C\bar A\)}
\end{frame}
